\section{The return in normalized touchdown coordinates}
\label{sec:detuned-return}
The contact theorem gives a release with normal velocity of order
$\delta^{3/2}$. The subsequent free gap is the integral of a normal
acceleration whose leading rate is $-c$, so its next zero has time and
scale of order $\delta^{3/4}$. To control the finer output fibres, we
cancel the leading free-flight terms before normalizing. We then
compare the compliant and reference returns at identical values of the
independent coordinates, and establish weighted state derivatives for
the exact reference return. These two estimates will control repeated
returns.

\subsection{Independent touchdown data and the physical inverse}
\label{sec:touchdown-coordinates}

Use the independent tuple \eqref{loc:tuple}, physical inverse
\eqref{loc:inverse} and input bounds \eqref{loc:domain}, writing
$c=c_\lambda$ and $\gamma=\gamma_{\rm cl}$.
The body coordinates are recovered by
\[
 \psi=y-1+\cos\phi_\delta,\qquad
 \dot\psi=v-\sin\phi_\delta\,\omega.
\]
For the inverse and comparison estimates below, the closing bound may
be any fixed $\gamma>0$. Constants depend only on the fixed input
bounds and compact model classes, with no positive lower bound on
$e$ or $r$.

\begin{lemma}[The dependent acceleration]
\label{ret:inverse}
On \eqref{loc:domain}, the physical inverse is $G+O(\delta)$ and
\begin{equation}
 \begin{aligned}
 b&=-c\delta+d_\lambda b_u(G)\delta^{4/3}\eta_U^r+b_u(G)r+O(\delta^2),\\
 \eta_B&=\frac{g_b}{c}
       \left(\eta_U^r+\frac{r}{d_\lambda\delta^{4/3}}\right)
       +O(\delta^{2/3}),\qquad g_b=-d_\lambda b_u(G).
 \end{aligned}
 \label{ret:inverse-expansion}
\end{equation}
In particular $|\eta_B|\le M_B$ for a common finite $M_B$ depending on
$M_{\rm in},A$ and the compact model classes, after reducing the common
scale cutoff. This bound is fixed before applying the contact theorem.
\end{lemma}
\begin{proof}
In the analytic chart $(p,y,u_0,v)$ the free rates at $G$ are
$(-a,0,d_\lambda,0)$. Therefore $F_{\rm free}b(G)=-c$ gives
\begin{equation}
 -a b_p(G)+d_\lambda b_u(G)=-c.\label{ret:rate-identity}
\end{equation}
Substitute \eqref{loc:tuple} with $u_0=u_r+r$ into the Taylor
expansion of $b$. Here $p,u_0=O(\delta)$, $v=O(\delta^2)$ and
$y=O(\delta^4)$; the remaining linear terms and quadratic remainder
are $O(\delta^2)$, uniformly on the fixed domain. Equation
\eqref{ret:rate-identity} yields the first formula, and division by
$-c\delta^{4/3}$ yields the second. The explicit inverse gives the
asserted bound on the mechanical state distance from $G$, together
with common bounds on all derivatives used.
\end{proof}

The exact reference $T_{0,\lambda}=(f_{0,\lambda},h_{0,\lambda})$
is defined by \eqref{loc:compression}--\eqref{loc:normal-reference}
and the subsequent release reconstruction and free flight in
Section~\ref{sec:local-results}. Write $Q=V(R_0)$ for its normal
release amplitude. The actual return constructed below is
$T_{\eps,r}=(f_{\eps,r},h_{\eps,r}^r)$.

\subsection{The first free event and normalization}
\label{sec:flight}

The normal fibre has denominator $D^{7/3}\asymp\delta^{7/4}$,
whereas its uncancelled velocity terms have size $\delta^{3/2}$.
We eliminate the release velocity with the touchdown equation before
dividing. For a free leg on its own elapsed release clock define
\begin{equation}
 \begin{gathered}
 D=-p(T)/a,\qquad J_U=u_0(T)+(d_\lambda/a)p(T),\\
 J_W=-v(T)-(c/3)D^2,\qquad
 A_p=\dot p+a,\quad B_b=\dot b+c.
 \end{gathered}\label{ret:critical-numerators}
\end{equation}
Integrating the free equations and using $y(T)=z(T)$ gives the exact
identities
\begin{align}
 D-T&=-p_L/a-a^{-1}\int_0^T A_p(t)\dd t,
 \label{ret:flight-offset}\\
 J_U&=u_{0,L}+(d_\lambda/a)p_L+
    \int_0^T\{a_T+(d_\lambda/a)\dot p\}\dd t,
 \label{ret:slip-cancellation}
\end{align}
For the normal numerator use the exact identity
\eqref{prep:touchdown-cancellation}, with $H_b=B_b$ and
$J_W=-v(T)-(c/3)D^2$. Thus the actual release velocity is eliminated
before either normalized output is formed.
The tuned output fibres are
\begin{equation}
 \eta_U^+=\frac{J_U}{d_\lambda D^{4/3}},\qquad
 \eta_W^+=\frac{J_W}{r_0D^{7/3}},\qquad
 \zeta^+=\frac{z_Le^{-T/\eps}}{\eps D^2}.
 \label{ret:normalization}
\end{equation}
For zero reference memory the last expression means $\zeta^+=0$.

\begin{lemma}[Paired free flights]
\label{lem:paired-free-normalization}
Consider either the actual and reference releases of
Theorem~\ref{con:contact-release} specialized by
\eqref{con:native-contact} at the same native input and offset,
or two exact releases with the same
$p_L=-a\delta$, $y_L=z_L=0$, $v_L=\delta^{3/2}Q$ and slips
$u_{0,L}=r,0$, where $Q$ lies in a fixed positive compact interval.
Put $L=\delta^{3/4}$ and set $\vartheta=e$ in the first case and
$\vartheta=r/\delta$ in the second. Both legs have first subsequent
touchdowns, positive gap between events, and
\begin{equation}
 \begin{gathered}
 T_i\asymp D_i\asymp L,\qquad
 g_i'(T_i),v_i(T_i)\le-\kappa_f\delta^{3/2},\\
 |\Delta T|+|\Delta D|\le C\vartheta L,\qquad
 |\log(D_1/D_0)|+|\xi_1^{0,+}-\xi_0^{0,+}|\le C\vartheta.
 \end{gathered}\label{ret:paired-values}
\end{equation}
On each full free prefix, $|X_i(\ell)-G|\le C(\delta+\ell)$ and
$T_i\le HL$. In the actual leg
\begin{equation}
 z(\ell)=z_Le^{-\ell/\eps}<0,\qquad
 |z(\ell)|/\eps\le C\delta^{3/2},\qquad
 |\zeta^+|\le C\exp\{-\kappa_f/(eL)\}\le CeL.
 \label{ret:free-memory}
\end{equation}
Here $\kappa_f>0$ is a common lower-bound constant. All output
coordinates in \eqref{ret:paired-values} use $u_0$.
\end{lemma}
\begin{proof}
The native contact estimate and the smooth inverse, respectively, give
\begin{equation}
 \begin{gathered}
 |X_{L,i}-G|\le C\delta,\quad b_{L,i}=O(\delta),\quad
 |\Delta p_L|+|\Delta u_0|+|\Delta b_L|\le C\vartheta\delta,\\
 |\Delta v_L|\le C\vartheta\delta^{3/2},\quad
 y_{L,0}=z_{L,0}=0,\quad y_{L,1}=z_{L,1},\quad
 |z_{L,1}|\le C\vartheta\delta^3.
 \end{gathered}\label{ret:paired-release-data}
\end{equation}
The actual release also has $|z_L|\le C\eps\delta^{3/2}$ and
$v_L+z_L/\eps\ge \kappa_{\rm op}\delta^{3/2}$. In the second pair,
$\Delta u_0=r$ and $\Delta b_L=O(r)$. Indeed, at fixed angle the
body velocity differences are exactly
\begin{equation}
 \Delta\omega=r/\cos\phi_\delta,\qquad
 \Delta\dot\psi=-r\tan\phi_\delta,
 \quad \Delta b_L=b_u(G)r+O(r\delta+r^2).
 \label{ret:body-shift}
\end{equation}
The two references have the same normal release coordinates, but
different mechanical velocities.

Stop the analytic free extensions at $HL$ or the chart boundary.
Their common smooth field gives
$|X_i(t)-G|\le C(\delta+t)$, excluding that boundary, and
$|\Delta X(t)|\le C\vartheta\delta$ by Gronwall's inequality.
Each residual $P\in\{A_p,B_b,a_T+(d_\lambda/a)\dot p\}$ vanishes at $G$,
so $|P_i(t)|\le C(\delta+t)$ and
$|\Delta P(t)|\le C\vartheta\delta$.
The exact gap is
\begin{equation}
 g_i(t)=v_{L,i}t+z_{L,i}(1-e^{-t/\eps})
                  +\int_0^t(t-s)b_i(s)\dd s.
 \label{ret:exact-free-gap}
\end{equation}
The memory term is absent when $z_L=0$. Since $z_L<0$, its
contribution is at least $(z_L/\eps)t$. The opening margin therefore
gives $g_i(t)\ge t(\kappa_f\delta^{3/2}-C\delta t-Ct^2)>0$ for
$0<t\le\eta L$, with fixed small $\eta>0$.
On $\eta\le x\le H$,
\[
 \delta^{-9/4}g_i(Lx)
 =Qx-cx^3/6+O_{C^1_x}(\delta^{1/4}+\vartheta).
\]
Here the differentiated memory is bounded using
$e^{-\eta/(eL)}\le Ce$. The positive interior portion and a common
transverse bracket about $\sqrt{6Q/c}$ identify the first zero.
The release and field differences give
\[
 |\Delta g(t)|\le C\vartheta
 (\delta^3+\delta^{3/2}t+\delta t^2)
 \le C\vartheta\delta^{9/4}.
\]
Dividing at the transverse bracket and then moving the event in $p$
proves the time, scale and logarithmic bounds. It also gives
$u_0(T_i)=d_\lambda T_i+O(\delta+T_i^2)>0$; subtracting $r=O(\delta^{4/3})$
preserves positive native slip.

For the normalized rows, substitute $t=T_iq$ in each residual integral.
The residuals there have size $O(L)$ and paired difference
$O(\vartheta L)$, including the changed argument. Equation
\eqref{ret:flight-offset} gives $H_i=D_i-T_i=O(\delta)$ and
$\Delta H=O(\vartheta\delta)$. Thus
$J_{U,i}=O(\delta)$ and $\Delta J_U=O(\vartheta\delta)$.
In \eqref{prep:touchdown-cancellation} keep
\[
 T_i^2-D_i^2=-H_i(2T_i+H_i)
\]
before dividing. The $b_LT$ and squared-scale terms have size
$O(\delta L)$ and difference $O(\vartheta\delta L)$.
The last integral has paired difference $O(\vartheta L^3)$,
and the memory term has size $O(\vartheta\delta^3/L)$.
Since $L^2\le\delta$, these give
$J_{W,i}=O(\delta L)$ and $\Delta J_W=O(\vartheta\delta L)$.
For either exponent $\alpha=4/3,7/3$ use
\begin{equation}
 \frac{J_1}{D_1^\alpha}-\frac{J_0}{D_0^\alpha}
 =\frac{\Delta J}{D_1^\alpha}
  +\frac{J_0}{D_0^\alpha}
       \{(D_1/D_0)^{-\alpha}-1\}.
 \label{ret:two-denominators}
\end{equation}
Because $D_i^{4/3}\asymp\delta$ and
$D_i^{7/3}\asymp\delta L$, this proves the fibre bound with their
respective denominators. The exact memory formula and
$D_i^2\asymp\delta^{3/2}$ prove \eqref{ret:free-memory}.
\end{proof}

\subsection{Weighted regularity of the exact limiting return}

Use the constants \eqref{loc:constants} and the weighted norm of
Section~\ref{sec:local-results}. Put
$B_\lambda=-ab_p(G)=-c+g_b$ and
$\mathscr D_\delta=\delta\partial_\delta$.
The reference derivatives are needed on a bounded open fibre domain,
including inputs outside the physical closing condition.

\begin{proposition}[Weighted limiting factorization]
\label{prop:limiting-factorization}
Let $\Omega\Subset\Omega_1$ be fixed bounded convex open sets containing
the independent input ball $|\xi|\le M_{\rm in}$ and the ball $\Kcal$
specified in Proposition~\ref{loc:return}.
On a common $(0,\delta_w]\times\overline\Omega$, the reference has
the $C^1$ extension, weighted bounds \eqref{loc:factorization} and joint
$\lambda$-continuity stated there. The memory component of $h_0$
vanishes. Moreover,
$f_0>\delta$ and
$\mathscr D_\delta\log f_0=3/4+O(\delta^{1/4})$.
\end{proposition}
On the auxiliary domain this is the algebraic extension of the signed
linear contact reference: compression is continued even if its force
is initially negative. No physical contact assertion is made there.
The actual return is used only under the closing condition in
\eqref{loc:domain}.
\begin{proof}
Lemma~\ref{con:reference-regularity} treats
the filter scale, frozen contact coefficients, initial normal
displacement and filter datum, slip and incoming normal velocity as
independent parameters. Restrict it to the native coefficient
$\pi=-a$ and write the remaining parameters as
$(\sigma,\widehat\delta,\alpha_U,\alpha_W,\beta_m)$, with
$h_0=h_{\widehat\delta}$, initial displacement and filter datum both
$\beta_m$, initial slip $d_\lambda(1+\alpha_U)$ and incoming speed
$r_0(1+\alpha_W)$. Its first-arrest and first-release maps are $C^1$
on a common neighborhood of zero in these independent parameters.
In particular the normal
release amplitude is $C^1$ there and has positive limit $q_\lambda$.
These derivative bounds include the independent initial displacement
and filter datum at $\sigma=0$; they do not use physical force
positivity for arbitrary auxiliary data.
Use the regularizing coordinate
\begin{equation}
 \rho=\delta^{1/12},\quad \sigma=\rho^6,\quad
 \widehat\delta=\rho^{12},\quad
 \alpha_i=\rho^4\eta_i,\quad \beta_m=\rho^6\zeta.
 \label{ret:regularizing-scale}
\end{equation}
The chain rule, including the term
$(\beta_m/2)\partial_{\beta_m}$ in $\mathscr D_\delta$, gives
$\|Q-q_\lambda\|_{1,*}\le C\rho^4$ and
$|D_\xi Q|\le C\rho^4$. Only the cutoff is reduced to cover the
fixed $\Omega_1$.

For an independent positive $Q_*$ near $q_\lambda$, reconstruct
$X_L=(-a\rho^{12},0,0,\rho^{18}Q_*)$ in the chart
$(p,y,u_0,v)$, centered at $G=0$. With $t=\rho^9s$ and
$X=\rho^9\widehat X$, its free equation becomes
$\widehat X_s=F_{\rm free}(\rho^9\widehat X)$ with initial state
$(-a\rho^3,0,0,\rho^9Q_*)$. This equation extends smoothly to
$\rho=0$. For every residual $P$ vanishing at $G$,
\[
 P(\rho^9\widehat X)=\rho^9
       \int_0^1DP(t\rho^9\widehat X)\widehat X\dd t.
\]
This divided form controls its remainders and derivatives. In particular
$b_L=\rho^{12}(B_\lambda+\rho^6B_1)$ and
\begin{equation}
 \rho^{-27}y(\rho^9s)
 =Q_*s+\tfrac12\rho^3(B_\lambda+\rho^6B_1)s^2
       -cs^3/6+\rho^9R_f,
 \label{ret:regular-free-gap}
\end{equation}
with smooth bounded remainders on fixed $s$ intervals. The positive
cubic root $\sqrt{6Q_*/c}$ has derivative denominator $-2Q_*$.
The opening interval and positive intervening gap in the proof of
Lemma~\ref{lem:paired-free-normalization} apply throughout this primitive domain;
thus the implicit root is the first positive root. It has the form
$S=\sqrt{6Q_*/c}+\rho^3S_1$. Integrating $p$ gives
\begin{equation}
 T_0=\rho^9S,\quad D_0=\rho^9\mathcal D,\quad
 \mathcal D=S+\rho^3+\rho^9H_1,
 \quad D_0-T_0=\rho^{12}(1+\rho^6H_1).
 \label{ret:regular-free-scale}
\end{equation}
After substituting the $C^1$ contact amplitude $Q(\rho,\xi)$,
$\mathcal D$ is bounded below and
$\|\mathcal D-K_\lambda\|_{1,*}\le C\rho^3$,
$|D_\xi\mathcal D|\le C\rho^4$.

The two normal cancellations remain visible in
\eqref{prep:touchdown-cancellation}: $D_0-T_0=\delta+O(\delta^{3/2})$
and $b_L=B_\lambda\delta+O(\delta^{3/2})$ give
\[
 J_W=(-B_\lambda/2-2c/3)\delta D_0
        +(B_\lambda/2+c/3)\delta^2+O(\delta^{9/4}).
\]
Using the divided remainders above makes this an identity with bounded
$C^1$ remainder, rather than differentiation of an order estimate.
Consequently there are bounded $C^1$ functions $U_1,W_1$ such that
\begin{equation}
 \begin{aligned}
 \eta_{U,0}^+&=(-1+\rho^6U_1/d_\lambda)\mathcal D^{-4/3},\\
 \eta_{W,0}^+&=-\varkappa_W\mathcal D^{-4/3}
       +A_1\rho^3\mathcal D^{-7/3}
       +\frac{\rho^6W_1}{r_0\mathcal D^{7/3}},
 \quad A_1=1+\frac{B_\lambda}{2r_0},\qquad \zeta_0^+=0.
 \end{aligned}\label{ret:regular-fibres}
\end{equation}
Their input-fibre dependence is only through $Q$. Product and chain
rules prove \eqref{loc:factorization}; the radial derivative of
$\log\rho^9$ is $3/4$.
\end{proof}

\subsection{A direct comparison at identical coordinate values}

The actual inverse \eqref{loc:inverse} uses $u_0=u_r+r$,
whereas $T_0(\delta,\xi^r)$ uses the tuned inverse at identical
values of the independent coordinates. These are generally different
physical touchdowns.
We pass between their returns through a reference with the actual
offset $r$, and recode the actual output after its scale is known.

\begin{theorem}[Value comparison at matched normalized coordinates]
\label{thm:detuned-value-bridge}
For sufficiently small inputs in \eqref{loc:domain}, the actual first
contact--flight return exists, has $D=f_{\eps,r}\asymp\delta^{3/4}$,
and satisfies \eqref{loc:value}.
The actual memory is the exact negative memory of
\eqref{ret:free-memory}. The first events, their transverse margins,
and all compact labels are retained.
\end{theorem}
\begin{proof}
Lemma~\ref{ret:inverse} places the actual input in a fixed finite
physical cone, so the native specialization
\eqref{con:native-contact} of Theorem~\ref{con:contact-release} applies.
Reconstruct
its reference release at $u_{0,L}=r$, with the same finite-$\sigma$
compression and inherited release, and call its exact free return
$R^r$. The first pair in Lemma~\ref{lem:paired-free-normalization} gives
\[
 |\log(f_{\eps,r}/f_{R^r})|
       +|h_{\eps,r}^{0}-h_{R^r}^{0}|\le Ce.
\]
For fixed $\delta$ the frozen angle and $h_\delta$ do not depend on
$r$. Hence $R^r$ and $T_0(\delta,\xi^r)$ share exactly the entire
compression solution, inherited release, first release root and
amplitude $Q$, not only their leading amplitudes. Their reconstructed
slips are $r$ and $0$. The second free pair, including the $O(r)$ body
and acceleration changes in \eqref{ret:body-shift}, gives
\[
 |\log(f_{R^r}/f_0)|+|h_{R^r}^{0}-h_0|\le Cr/\delta.
\]
In particular $\Delta b_LT=O(rL)$ is retained in the normal numerator
before division by $D^{7/3}\asymp\delta L$.
Finally recode only the actual output:
\begin{equation}
 h_{\eps,r}^r=h_{\eps,r}^0-
               \frac{r}{d_\lambda D^{4/3}}(0,1,0).
 \label{eq:dva-output-recoding}
\end{equation}
Its cost is $Cr/\delta$ because the paired free flights have already
established $D^{4/3}\asymp\delta$. Summing these three comparisons
proves the theorem.
\end{proof}

For complete physical-state comparisons, recover the acceleration
from the same analytic inverse chart. Define
\[
 \mathcal B_{\eps,r}(s,\xi)=
 -\frac{b(-as,\eps s^2\zeta,d_\lambda s+d_\lambda s^{4/3}\eta_U+r,
                 -r_0s^2-r_0s^{7/3}\eta_W)+cs}{cs^{4/3}}.
\]
Taylor's formula and \eqref{ret:rate-identity} give the first four
bounds below at $r=0$, together with the offset comparison in the
last bound:
\begin{equation}
 |\partial_{\eta_U}\mathcal B|\le C,\quad
 |\partial_{\eta_W}\mathcal B|\le Cs,\quad
 |\partial_\zeta\mathcal B|\le C\eps s^{2/3},\quad
 |s\partial_s\mathcal B|\le Cs^{2/3},\quad
 |\mathcal B_{\eps,r}-\mathcal B_{\eps,0}|\le Cr/s^{4/3}.
 \label{ret:dependent-Lipschitz}
\end{equation}
The radial derivative keeps $\eps,r$ fixed; its leading cancellation
is exactly \eqref{ret:rate-identity}. Integrating these estimates on
the logarithmic scale interval between the two output scales and on
the bounded fibre segment gives
$|\eta_B^+-\eta_{B,0}^+|\le C(e+r/\delta)$.
The memory component of $h_0$ vanishes, so inserting the same
$\eps$ into its inverse is exact. This recovery uses both output scales; the
independent state remains $(\log\delta,\zeta,\eta_U^r,\eta_W)$.

\subsection{Admissibility of the returned state}
\label{sec:detuned-output-admission}

For the already fixed closed ball $\Kcal$, put
$d_\Kcal=\inf_\lambda\operatorname{dist}
(\xi_{*,\lambda},\partial\Kcal)>0$, and retain
$\gamma=\gamma_{\rm cl}$ and $M_{\rm in}$ from \eqref{loc:domain}.
The ball and independent bound are fixed once and for all before
the return estimates; the latter contains the second-touchdown tuples
of \eqref{prep:complete-second-cone} and $\Kcal$ in its interior.
The dependent acceleration is recovered from the physical inverse.

\begin{theorem}[Admissibility of the returned state]
\label{thm:detuned-output-admission}
With these choices, the common radii, output bounds
\eqref{loc:admission}, signs and inverse branch in
Proposition~\ref{loc:return} hold, including its inner and outer
domain qualifications and order of radius choices. The closing
margin is $r_0W^+-\zeta^+>\gamma$, and the dependent acceleration
is uniformly bounded as in Lemma~\ref{ret:inverse}.
\end{theorem}
\begin{proof}
The compact classes, independent bound $M_{\rm in}$,
$\Kcal,\Omega,\Omega_1,A,C_0,\gamma$ and inverse neighborhood are fixed.
Lemma~\ref{ret:inverse} gives a finite $M_B$ depending on
$M_{\rm in},A$ and the compact model classes. Choose a physical cone
containing these independent and dependent bounds, and then obtain
the contact, free and reference constants and a common
$\delta_{\rm call}$. Choose any sufficiently small
$\delta_{\rm g}<\delta_{\rm call}$ within the limiting domain, with
$\delta_{\rm g}^{1/3}\sup_{\Kcal}(|\eta_U|+|\eta_W|)\le1/4$.
Theorem~\ref{thm:detuned-value-bridge} and
Proposition~\ref{prop:limiting-factorization} prove the estimate in
\eqref{loc:admission}. Its error tends to zero uniformly because
$e\le C_0\delta^{1/2}$ and $r/\delta\le A\delta^{1/3}$.
Choose $\delta_{\rm reset}$ last, making this error less than
$d_\Kcal/2$ and ensuring, for the common scale bounds
$c_D\delta^{3/4}\le D\le C_D\delta^{3/4}$,
$c_D\delta_{\rm reset}^{-1/4}>1$ and
$C_D\delta_{\rm reset}^{3/4}<\delta_{\rm g}$.
Then $D>\delta$, so the compliance and offset conditions propagate.
The ratio bounds give $U^{r,+},W^+\ge3/4$ and hence the velocity
signs. Since $\zeta^+<0$,
$r_0W^+-\zeta^+\ge(3/4)\min r_0>\gamma$.
The same inverse applies, and Lemma~\ref{ret:inverse} recovers its
dependent coordinate on the enlarged output scale. The label and
physical parameters were fixed throughout.
\end{proof}

The two smallness conditions in \eqref{loc:domain} have different
roles: $\eps\le C_0\delta^2$ controls compliance, while
$r\le A\delta^{4/3}$ bounds the dependent acceleration in the
physical inverse. The first contact uses the broader range
$\delta=\sqrt\eps$, $r=k\delta$ of the contact theorem; the two
preparatory flights bring the same trajectory into \eqref{loc:domain}.
